% GENERATED FILE. Edit canonical lessons, then run npm run generate:books.
% canonical-source-hash: fnv1a64:0e5f81b490f2a9a0
% canonical-lessons: FA-C381-sarparast, FA-C381-moaven, FA-C381-karshenas, FA-C381-motakhasses, FA-C381-vazife

\chapter{Roles at work}
\label{ch:fa-roles-at-work}
\hlchaptermodality{381}

\begin{chapteropening}
\textbf{By the end of this chapter:} \emph{I can name roles at work, such as a supervisor, a deputy, an expert and a specialist, and talk about my duties.}

Complete the last lesson of chapter 381: I can name roles at work, such as a supervisor, a deputy, an expert and a specialist, and talk about my duties.
\end{chapteropening}

\section[\texorpdfstring{\fa{سرپرست} (sarparast)}{سرپرست (sarparast)}]{\fa{سرپرست} (sarparast) — a supervisor; a guardian}

\label{lesson:FA-C381-sarparast}



\begin{quote}
Before the new one: say the Persian for a trainee, then the Persian for an internship.
\end{quote}

\subsection*{You'll want to know: \fa{سرپرست}}
\textbf{\fa{سرپرست}} (\emph{sarparast}) — \textquotedblleft{}a supervisor; a guardian\textquotedblright{}.

\begin{quote}
\fa{سرپرست} \fa{ما} \fa{هر} \fa{صبح} \fa{کارها} \fa{را} \fa{تقسیم} \fa{میکند}.
\end{quote}

\begin{quote}
\emph{sarparast-e mâ har sobh kârhâ râ taqsim mikonad.}
\end{quote}

\begin{quote}
Our supervisor shares out the jobs every morning.
\end{quote}

\emph{Sar} \textquotedblleft{}head\textquotedblright{} + \emph{parast}, from \emph{parastidan} \textquotedblleft{}to serve, to tend\textquotedblright{}. Also the guardian of a child.

\begin{cousinweb}
\textbf{\fa{سرپرست}} is a native Persian compound built from \emph{sar} \textquotedblleft{}head\textquotedblright{} + \emph{parast}, which EtymDB traces to a Middle Persian verb \textquotedblleft{}to serve, to minister\textquotedblright{}. A supervisor is the one who looks after the head of the work.
\end{cousinweb}

\begin{grammarlens}[title={What you've built}]
One more word for telling, explaining and describing at B1.
\end{grammarlens}

\subsection*{Your turn}
\begin{itemize}
\raggedright
  \item \emph{Say it:} \emph{sarparast}
  \item \emph{Say it:} \emph{sarparast}, once more
  \item \emph{Say it:} a sentence of your own with \emph{sarparast}
  \item \emph{Recall:} say the Persian for a trainee, then the Persian for an internship, then say \emph{sarparast} again
\end{itemize}

\begin{tcolorbox}[breakable,colback=teal!4,colframe=teal!35!black,title={Before you move on}]
What does \fa{سرپرست} mean? (\textquotedblleft{}A supervisor; a guardian\textquotedblright{}.) Say it once more.
\end{tcolorbox}
\section[\texorpdfstring{\fa{معاون} (mo'âven)}{معاون (mo'âven)}]{\fa{معاون} (mo'âven) — a deputy, an assistant manager}

\label{lesson:FA-C381-moaven}



\begin{quote}
Before the new one: say the Persian for an internship, then the Persian for a supervisor.
\end{quote}

\subsection*{You'll want to know: \fa{معاون}}
\textbf{\fa{معاون}} (\emph{mo'âven}) — \textquotedblleft{}a deputy, an assistant manager\textquotedblright{}.

\begin{quote}
\fa{معاون} \fa{مدرسه} \fa{با} \fa{والدین} \fa{صحبت} \fa{کرد}.
\end{quote}

\begin{quote}
\emph{mo'âven-e madrese bâ vâledeyn sohbat kard.}
\end{quote}

\begin{quote}
The deputy head of the school talked with the parents.
\end{quote}

Plural \emph{mo'âvenân}. \emph{Mo'âven-e vazir} — a deputy minister.

\begin{cousinweb}
\textbf{\fa{معاون}} is usually traced to Arabic \emph{mu'āwin} \textquotedblleft{}helper, assistant\textquotedblright{}: the person who helps the head.
\end{cousinweb}

\begin{grammarlens}[title={What you've built}]
One more word for telling, explaining and describing at B1.
\end{grammarlens}

\subsection*{Your turn}
\begin{itemize}
\raggedright
  \item \emph{Say it:} \emph{mo'âven}
  \item \emph{Say it:} \emph{mo'âven}, once more
  \item \emph{Say it:} a sentence of your own with \emph{mo'âven}
  \item \emph{Recall:} say the Persian for an internship, then the Persian for a supervisor, then say \emph{mo'âven} again
\end{itemize}

\begin{tcolorbox}[breakable,colback=teal!4,colframe=teal!35!black,title={Before you move on}]
What does \fa{معاون} mean? (\textquotedblleft{}A deputy, an assistant manager\textquotedblright{}.) Say it once more.
\end{tcolorbox}
\section[\texorpdfstring{\fa{کارشناس} (kârshenâs)}{کارشناس (kârshenâs)}]{\fa{کارشناس} (kârshenâs) — an expert, an analyst}

\label{lesson:FA-C381-karshenas}



\begin{quote}
Before the new one: say the Persian for a supervisor, then the Persian for a deputy.
\end{quote}

\subsection*{You'll want to know: \fa{کارشناس}}
\textbf{\fa{کارشناس}} (\emph{kârshenâs}) — \textquotedblleft{}an expert, an analyst\textquotedblright{}.

\begin{quote}
\fa{یک} \fa{کارشناس} \fa{خانه} \fa{را} \fa{دید} \fa{و} \fa{قیمت} \fa{گذاشت}.
\end{quote}

\begin{quote}
\emph{yek kârshenâs khâne râ did va qeymat gozâsht.}
\end{quote}

\begin{quote}
An expert looked at the house and set a price.
\end{quote}

\emph{Kâr} \textquotedblleft{}work\textquotedblright{} + \emph{shenâs} \textquotedblleft{}knowing\textquotedblright{}. \emph{Kârshenâsi} is also a bachelor's degree.

\begin{cousinweb}
\textbf{\fa{کارشناس}} is native Persian, built from \emph{kâr} \textquotedblleft{}work\textquotedblright{} + \emph{shenâs-}, the present stem of \emph{shenâkhtan} \textquotedblleft{}to know\textquotedblright{}, which EtymDB traces through Old Persian to PIE \emph{ǵneh\textsubscript{3}-} \textquotedblleft{}to know\textquotedblright{}.
\end{cousinweb}

\begin{grammarlens}[title={What you've built}]
One more word for telling, explaining and describing at B1.
\end{grammarlens}

\subsection*{Your turn}
\begin{itemize}
\raggedright
  \item \emph{Say it:} \emph{kârshenâs}
  \item \emph{Say it:} \emph{kârshenâs}, once more
  \item \emph{Say it:} a sentence of your own with \emph{kârshenâs}
  \item \emph{Recall:} say the Persian for a supervisor, then the Persian for a deputy, then say \emph{kârshenâs} again
\end{itemize}

\begin{tcolorbox}[breakable,colback=teal!4,colframe=teal!35!black,title={Before you move on}]
What does \fa{کارشناس} mean? (\textquotedblleft{}An expert, an analyst\textquotedblright{}.) Say it once more.
\end{tcolorbox}
\section[\texorpdfstring{\fa{متخصص} (motakhasses)}{متخصص (motakhasses)}]{\fa{متخصص} (motakhasses) — a specialist (doctor), an expert}

\label{lesson:FA-C381-motakhasses}



\begin{quote}
Before the new one: say the Persian for a deputy, then the Persian for an expert.
\end{quote}

\subsection*{You'll want to know: \fa{متخصص}}
\textbf{\fa{متخصص}} (\emph{motakhasses}) — \textquotedblleft{}a specialist (doctor), an expert\textquotedblright{}.

\begin{quote}
\fa{دکتر} \fa{مرا} \fa{پیش} \fa{یک} \fa{متخصص} \fa{قلب} \fa{فرستاد}.
\end{quote}

\begin{quote}
\emph{doktor marâ pish-e yek motakhasses-e qalb ferestâd.}
\end{quote}

\begin{quote}
The doctor sent me to a heart specialist.
\end{quote}

Above all for doctors and skilled trades. Plural \emph{motakhassesân}.

\begin{cousinweb}
\textbf{\fa{متخصص}} is usually traced to Arabic \emph{mutakhaṣṣiṣ} \textquotedblleft{}specialised\textquotedblright{}, from the same Arabic family as \emph{khosusi} \textquotedblleft{}particular, private\textquotedblright{}.
\end{cousinweb}

\begin{grammarlens}[title={What you've built}]
One more word for telling, explaining and describing at B1.
\end{grammarlens}

\subsection*{Your turn}
\begin{itemize}
\raggedright
  \item \emph{Say it:} \emph{motakhasses}
  \item \emph{Say it:} \emph{motakhasses}, once more
  \item \emph{Say it:} a sentence of your own with \emph{motakhasses}
  \item \emph{Recall:} say the Persian for a deputy, then the Persian for an expert, then say \emph{motakhasses} again
\end{itemize}

\begin{tcolorbox}[breakable,colback=teal!4,colframe=teal!35!black,title={Before you move on}]
What does \fa{متخصص} mean? (\textquotedblleft{}A specialist (doctor), an expert\textquotedblright{}.) Say it once more.
\end{tcolorbox}
\section[\texorpdfstring{\fa{وظیفه} (vazife)}{وظیفه (vazife)}]{\fa{وظیفه} (vazife) — a duty, a task}

\label{lesson:FA-C381-vazife}



\begin{quote}
Before the new one: say the Persian for an expert, then the Persian for a specialist.
\end{quote}

\subsection*{You'll want to know: \fa{وظیفه}}
\textbf{\fa{وظیفه}} (\emph{vazife}) — \textquotedblleft{}a duty, a task\textquotedblright{}.

\begin{quote}
\fa{کمک} \fa{به} \fa{مردم} \fa{وظیفه} \fa{ماست}.
\end{quote}

\begin{quote}
\emph{komak be mardom vazife-ye mâst.}
\end{quote}

\begin{quote}
Helping people is our duty.
\end{quote}

Plural \emph{vazâyef}, an Arabic broken plural. \emph{Vazife dâshtan} — to be obliged to.

\begin{cousinweb}
\textbf{\fa{وظیفه}} was borrowed from Arabic \emph{waẓīfa} \textquotedblleft{}task, duty, office\textquotedblright{} (EtymDB).
\end{cousinweb}

\begin{grammarlens}[title={What you've built}]
One more word for telling, explaining and describing at B1.
\end{grammarlens}

\subsection*{Your turn}
\begin{itemize}
\raggedright
  \item \emph{Say it:} \emph{vazife}
  \item \emph{Say it:} \emph{vazife}, once more
  \item \emph{Say it:} a sentence of your own with \emph{vazife}
  \item \emph{Recall:} say the Persian for an expert, then the Persian for a specialist, then say \emph{vazife} again
\end{itemize}

\begin{tcolorbox}[breakable,colback=teal!4,colframe=teal!35!black,title={Before you move on}]
What does \fa{وظیفه} mean? (\textquotedblleft{}A duty, a task\textquotedblright{}.) Say it once more.
\end{tcolorbox}
